% discussion.tex — Full proposal draft, Week 8
% Synthesized from Week 4 evaluation and Week 7 limitations, with new framing.

\section{Discussion}

\subsection{What the Preliminary Results Establish}

The NNW training phase establishes three things that are necessary before the cross-task analyses can proceed. First, the aPCR model significantly predicts 56 of 122 brain regions during naturalistic movie-watching, demonstrating that whole-brain fMRI dynamics can be recovered from surface fNIRS signals in a dataset where fNIRS and fMRI were recorded from entirely different people. Second, the predicted connectivity structure closely matches the observed connectivity structure ($r = 0.617$), which means the model is not just predicting isolated regions but preserving the large-scale organizational logic of the brain, including the antagonistic relationship between task-positive and default-mode networks. Third, the networks most relevant to the individual-differences hypothesis (DAN, CONT, DMN) are among the best-predicted, with 29 of 56 significant ROIs falling in these three systems. This pattern is consistent with the fact that these networks are the most consistently engaged by the narrative structure of a film \citep{hasson2010} and are therefore the networks where the cross-participant, cross-modality alignment is strongest.

\subsection{Interpreting the Network Ranking Discrepancy}

One finding that requires careful interpretation is that the network-level prediction profile in the primary dataset differs from both the \citet{gao2025} framework and an independent replication dataset collected in the same lab. In the replication dataset, DMN and CONT lead in prediction accuracy, consistent with Gao et al.'s results using BBC \textit{Sherlock}. In the primary dataset, VIS and DAN lead, with DMN and CONT performing well but not dominating. Both datasets produce significant whole-brain predictions with meaningful ISFC, but they emphasize different networks.

A plausible explanation involves session ordering and time-on-task effects. In the primary dataset, participants watched NNW at the end of a two-hour session after completing all cognitive tasks. A well-established literature on time-on-task effects shows that prolonged cognitive effort alters the dynamic balance between task-positive and default-mode networks: sustained attention depletes executive resources and reduces the efficiency of DMN suppression and reactivation \citep{sonugabarke2007}. If fatigue selectively dampens narrative-driven DMN engagement with the film (because tired participants track the story less deeply) while leaving stimulus-driven visual and attentional responses relatively intact (because bottom-up processing of visual action sequences requires less effortful engagement), this would shift the network profile toward VIS and DAN without degrading overall prediction accuracy. The replication dataset, which used a different session structure, would not show this shift, explaining why it reproduces the Gao et al.\ pattern.

This interpretation has a testable implication. Future data collection will place the movie-watching block first in the session, before any cognitive tasks, and will include self-report measures of engagement and cap discomfort. If the movie-first ordering restores the DMN/CONT-dominant pattern, that would confirm that session ordering modulates which networks the model predicts best. This also has practical consequences for the individual differences analyses in Phase 5: the DMN-RTV hypothesis depends on the model's ability to predict DMN dynamics during cognitive tasks. If DMN prediction is attenuated in the primary dataset due to training-phase fatigue effects, the DMN-RTV correlation may be weaker than expected, making the movie-first design change important for maximizing the sensitivity of the individual differences test.

\subsection{Why the Cross-Task Test Is Meaningful}

The most principled aspect of this design is using movie-watching as the training stimulus rather than a structured task. The whole analysis depends on having temporally consistent brain responses across both fNIRS and fMRI participants, and naturalistic stimuli are what produce that consistency \citep{hasson2010, nastase2019}. Structured cognitive tasks could not serve this role because individual differences in performance, strategy, and arousal make brain responses too person-specific to train a cross-dataset model. So the design is not using movie-watching out of convenience; it is using it because it is the paradigm that best supports cross-person, cross-modality model training. The whole logic of the project is to calibrate on shared signal, then test on individual variation.

The frozen-model design also directly addresses what would otherwise be an insurmountable practical constraint: there is no cognitive-task fMRI dataset from the same people who did the fNIRS cognitive tasks. Rather than treating this as a problem, the design frames it as the test itself: does a model calibrated on passive movie-watching transfer to structured cognitive tasks? The NNW results show the model produces coherent predictions in the training domain, which justifies testing cross-task transfer.

\subsection{Limitations}

Several limitations should be noted.

The main limitation is spatial coverage. The fNIRS cap covers prefrontal and lateral parietal cortex but has no channels over medial prefrontal cortex, primary visual cortex, or subcortical areas. The model's ability to predict those regions during cognitive tasks rests entirely on statistical covariation patterns learned from the movie. The NNW results show this can work (VIS and SM are significantly predicted despite no direct coverage), but the mechanism is likely shared temporal dynamics during the movie, not true spatial fidelity. This is an important caveat for any ROI-level interpretation of cognitive task predictions.

Sample size is a concern. N\,=\,28 for the NNW model is enough to test the pipeline, but the individual differences regression in Phase 5 will struggle to detect effects smaller than $r \approx 0.40$ at this sample size. The plan is to keep recruiting toward N\,=\,40, and to report effect sizes and confidence intervals rather than relying solely on significance, so the results are interpretable regardless of where the final sample lands.

There is no cognitive-task fMRI ground truth to directly validate the predictions against. The validation strategy, checking ISFC, comparing against the n-back literature, and correlating with behavior, is indirect but informative. Each of these provides a different kind of evidence that the predictions are meaningful, and together they make a reasonable case even without matched fMRI data. External validation against large public fMRI datasets (HCP n-back, N\,=\,1,200; DMCC AX-CPT, N\,=\,55) remains an option if the within-sample correlations are too underpowered.

Finally, the behavioral data quality issues described in the Methods (duplicate files, accuracy coding verification, CPT metric recomputation) must be resolved before the individual differences analyses can be run. These are bookkeeping problems, not conceptual ones, but they need to be completed before Phase 5 proceeds.

\subsection{Why This Is the Right Approach}

The obvious alternative would be to collect fMRI and fNIRS simultaneously during cognitive tasks from the same participants and train the model directly on that data. This would solve the spatial validation problem and eliminate the frozen-model assumption. But it also requires every participant to go in the scanner, which is the exact problem the project is trying to solve. If scanner access is required, there is no scalable, equitable version of this tool.

The naturalistic training design trades direct validation for accessibility, and that tradeoff is the whole point of the project. The NNW results show the model is not noise; it captures real structure during the training condition. The cognitive task analyses will tell us how much of that structure persists when the context changes. Whatever the answer is, it moves the field forward on a question that matters both scientifically and practically.

Even a null result has scientific value. If the frozen model fails to predict cognitive-task brain patterns or fails to correlate with attention performance, that is a meaningful boundary condition on the approach, not just a failed experiment. It would tell us something specific about when naturalistic calibration does and does not generalize, which is useful for the broader goal of portable neurocognitive assessment \citep{farah2006, noble2015}.

\subsection{Broader Implications}

If the cross-task predictions succeed, the practical implications extend well beyond this study. A portable fNIRS session lasting under 15 minutes (the NNW clip), requiring no scanner, no special infrastructure, and no trained MRI technician, would be sufficient to calibrate a model that then provides whole-brain-equivalent neural measures during any subsequent cognitive assessment. This opens the possibility of neurocognitive research in schools, community clinics, field settings, and populations that have been structurally excluded from scanner-based science. The equity argument runs through every design choice in this project, from the use of a shared movie stimulus that requires no literacy or cultural familiarity, to the focus on individual differences in attention that track socioeconomic and clinical disparities, to the explicit goal of building a tool that works without the infrastructure that currently gatekeeps access to brain science \citep{burns2019, farah2006}.
